% !PRINTED 191
\weberpage{191}
%<<CONT
und wir müssen also die Determinante $\alpha\delta - \beta\gamma = \pm 1$
setzen. Berechnet man die Determinante $|A'|$, so ergiebt sich
dafür $(\alpha\delta - \beta\gamma)^3$, so dass also, wenn wir nur
die eigentlichen orthogonalen Substitutionen berücksichtigen,
\weq{13}{\alpha\delta - \beta\gamma = 1}
setzen müssen, während
\weq{14}{\alpha\delta - \beta\gamma = -1}
den uneigentlichen orthogonalen Substitutionen entspricht.
Dadurch ist die Substitution $A'$ völlig bestimmt, und ist
zurückgeführt auf die binäre Substitution
\weq{15}{(\eta_1, \eta_2) = \begin{pmatrix} \alpha, & \beta \\ \gamma, & \delta \end{pmatrix} (\xi_1, \xi_2),\quad \alpha\delta - \beta\gamma = \pm 1,}
oder, wenn $\eta_1 : \eta_2 = \eta$, $\xi_1 : \xi_2 = \xi$ gesetzt
wird, auf die lineare gebrochene Substitution
\weq{16}{\eta = \dfrac{\alpha\xi + \beta}{\gamma\xi + \delta}.}
Hierzu ist aber nun noch Folgendes zu bemerken. Die beiden
Substitutionen $A, A'$ sind Transformationen von einander, und
entsprechen sich also gegenseitig eindeutig.

Durch $A'$ sind aber die Zahlen $\alpha, \beta, \gamma, \delta$
nur bis auf das gemeinsame Vorzeichen bestimmt. Wenn wir also die
lineare Substitution
\weq{17}{\begin{pmatrix} \alpha, & \beta \\ \gamma, & \delta \end{pmatrix}}
mit einer der beiden Bedingungen (13) und (14) betrachten, so ist
zwar hierdurch die Substitution $A'$ und daher auch $A$ eindeutig
bestimmt; aber umgekehrt entsprechen jedem $A'$ zwei
Substitutionen der Form (17):
\[
\begin{pmatrix} \alpha, & \beta \\ \gamma, & \delta \end{pmatrix}
\quad\text{und}\quad
\begin{pmatrix} -\alpha, & -\beta \\ -\gamma, & -\delta \end{pmatrix}.
\]
Wir erhalten aber wieder ein eindeutiges Entsprechen, wenn wir
diese beiden Substitutionen zu einem gemeinsamen Begriffe, den
wir mit $A''$ bezeichnen, zusammenfassen.

Was nun endlich die Substitution (16) betrifft, die wir mit
\weq{18}{A''' = \begin{pmatrix} \alpha, & \beta \\ \gamma, & \delta \end{pmatrix}}
bezeichnen, und nach \S.\,38 als Substitution der Verhältnisse
auffassen, so können wir einen gemeinsamen Factor von
$\alpha, \beta, \gamma, \delta$
%CONT>>
